\begin{frame}
\frametitle{Vliv velikosti a organizace cache na počty výpadků}

{
\centering

\includegraphics[width=0.65\linewidth]{generated/memory/cache/cache-size-vs-ways.pdf}

}

\begin{itemize}
\item miss rate není vlastnost, parameter cache
\item miss rate není vlastnost programu
\end{itemize}

Jedná se o kombinaci jak algoritmů v programu, tak parametrů cache a často i zpracovávaných dat.

\end{frame}

\begin{frame}
\frametitle{Vliv velikosti bloku a celé cache na počty výpadků}

Miss rate versus block size and cache size on SPEC92 benchmark

\vskip 2mm

{
\centering

\includegraphics[width=0.72\linewidth]{generated/memory/cache/cache-size-vs-block.pdf}

}

\vskip 2mm

{\tiny Zdroj: Hennessy and Patterson: Computer Architecture: A Quantitative Approach 3rd ed., Morgan Kaufmann, 2003}

Zvětšení velikosti bloku napomáhá načtení sousedních dat, která mohou být často následně potřeba.
Ale pokud tomu tak není, tak se zvyšuje cache miss penalty a zároveň dochází častěji ke kolizím
a obsazení cache nepotřebnými daty.

\end{frame}
